\section{Spherical oscillatory integrals and distance shells}
\label{osc:section}

This Section proves a Fourier estimate for weighted distance measures
built from separated class laws and endpoint masks.  Its output is a
spherical angular energy with the same endpoint data.  The proof uses
global stationary-phase operators on the sphere, so their derivatives
can be distributed between the two endpoint symbols without creating a
coefficient that depends jointly on the endpoints.  Section~\ref{rec:section}
will apply this estimate to high frequencies in a positive count of pairs
with nearly equal distances.

Throughout this Section, spatial and frequency depths are integers in
$[0,N]$.  In an angular energy they satisfy $0\le a\le v\le N$.
The largest nominal length of an empty descriptor list is zero.
Write $n_*=d-1$.

\subsection{Separated cells and endpoint symbols}

Fix a power of two $K_d$ with
$1000(d+1)<K_d<4000(d+1)$.  Two depth-$a$ cells $X,Y$ are
called \emph{separated} if
\begin{equation}\label{osc:separation}
 K_d2^{-a}/8\le |c_X-c_Y|\le 4K_d2^{-a}.
\end{equation}
For $x\in X$, $y\in Y$ we then have $D(x,y)\asymp_d2^{-a}$, and their
directions $e(x,y)=(x-y)/D(x,y)$ lie in a cone of small fixed aperture.
Choose $0\le\psi\le1$ smooth on $S^{d-1}$, identically one on a fixed
neighborhood of these directions and zero on a fixed neighborhood of
their negatives.  Its derivatives of every fixed order are uniformly
bounded.  We use the same convention for the fixed separated outer
supports, with $a=0$; constants there may depend on the initial geometry.

All angular derivatives below can be taken in the fixed smooth tangent
fields
\[
 \Omega_{uv}=w_u\partial_{w_v}-w_v\partial_{w_u},
 \qquad 1\le u<v\le d.
\]
These fields span every tangent space of the sphere.  A word of length
$q$ in these fields is denoted by $\Omega^\alpha$, $|\alpha|=q$; the empty
word is the identity.  Recall that a descriptor $\mathfrak t$ of nominal
length $l_{\mathfrak t}$ has angular scale
\[
 D_{l_{\mathfrak t}}=\max(1,2^{l_{\mathfrak t}}R^{-50\sigma}).
\]

\begin{definition}[Admissible symbols]\label{osc:symbols}
Fix an indexed descriptor list
$\mathcal A=(\mathfrak t_\nu)_{\nu\in\mathcal I}$ with
$\#\mathcal I\le N^{C_{\mathcal A}}$, a smoothing index
$i\in\{0,\ldots,J_0-1\}$, and a nonnegative integer derivative budget
$B$.  The exponent $C_{\mathcal A}$ and budget $B$ are fixed independently
of $N$ and are common to the symbol family under consideration.  At each
angular order under consideration, also fix a common numerical bound
$Q_j$ independent of $N$.  An admissible endpoint symbol has the form
\begin{equation}\label{osc:symbol-form}
 b_x(w)=f(x)q(w)
       \prod_{\nu\in\mathcal I}
       D_{l_{\mathfrak t_\nu}}^{-|\alpha_\nu|}
       \Omega^{\alpha_\nu}w_{i,\mathfrak t_\nu}(x,w),
 \qquad \sum_{\nu\in\mathcal I}|\alpha_\nu|\le B.
\end{equation}
The index $\nu$ distinguishes occurrences, so repeated descriptors may
receive different derivative words.
Here $f$ is an arbitrary measurable spatial function with $|f|\le1$,
and $q$ is a smooth angular function satisfying
$\|q\|_{C^j(S^{d-1})}\le Q_j$ at every order under consideration.
The constants in estimates may depend on $C_{\mathcal A}$, $B$ and the
fixed $Q_j$.
Fixed scalar factors can be absorbed in those estimates.  Empty lists
and empty products are allowed, as are complex conjugates of such
symbols.  Symbols do not depend on the radial variable $r$.
\end{definition}

No spatial regularity is required in Definition~\ref{osc:symbols}.
In particular, $f$ can impose an arbitrary Borel restriction inside a
cell.  Every differentiated mask has the same angular support as, or a
smaller support than, the original mask.  Thus every descriptor present
in a symbol still imposes its level-$(i+1)$ good condition.

\begin{lemma}[Derivative bounds and reached anchors]
\label{osc:symbol-bounds}
Let $L$ be the largest nominal length in the symbol's list, or in all
lists for a product of symbols, with $L=0$ when those lists are empty, and put
\begin{equation}\label{osc:derivative-scale}
 D_*=R^\sigma\max(1,2^L R^{-50\sigma}).
\end{equation}
For every fixed $j$, an admissible symbol, or a product of a fixed number
of endpoint symbols, obeys
\begin{equation}\label{osc:angular-derivatives}
 \sup_x\|b_x\|_{C^j(S^{d-1})}\le C_{j,B}D_*^j
\end{equation}
for sufficiently large $N$.  The constant may depend on the fixed
list-size exponents, parameters, derivative budgets and common angular
bounds, but not on the spatial multipliers.  On a depth-$a$ cell, the product of all mask factors
whose anchors have depth at most $a$ is independent of the spatial
integration variable and has uniformly bounded modulus.  We call these
anchors \emph{reached} at depth $a$.
\end{lemma}

\begin{proof}
Apply Lemma~\ref{test:product-symbols} to each endpoint symbol, using
the polynomial list bound, the prescribed derivative budget and the common
bounds for $q$.  A fixed product of endpoint symbols then satisfies the
same estimate by the finite product rule, even if their smoothing indices
differ.  For the last
assertion, an anchor cell at depth at most $a$ contains the depth-$a$
cell, and its mask is constant in the spatial variable on that cell.
Angular differentiation preserves this constancy.  The product has
bounded modulus because only boundedly many factors were prescribed
derivatives and all other masks have modulus at most one.
\end{proof}

\subsection{General angular energies}

For occupied depth-$a$ cells of one or two classes, with specified
admissible endpoint symbols, define
\begin{equation}\label{osc:transforms}
 U_X(r,w)=\frac1{M_1(a)}\int_X e^{irw\cdot x}b_x(w)\,d\rho_1(x),
 \qquad
 V_Y(r,w)=\frac1{M_2(a)}\int_Y e^{irw\cdot y}c_y(w)\,d\rho_2(y).
\end{equation}
These transforms are normalized by the profile model masses, not by the
actual cell masses.  Before the endpoint factors are inserted, the first
underlying measure has mass $\rho_1(X)/M_1(a)=R^{O_d(\sigma)}$, and
similarly on the second side.  It need not be a conditional probability.
Fix $\Psi\ge0$ in $C_c^\infty((1/4,4))$ with $\Psi\ge1$ on $[1/2,2]$.
Surface measure on the sphere is denoted by $dw$.

\begin{definition}[General angular energies]\label{osc:states}
For $0\le a\le v\le N$, the independent angular energy is
\begin{equation}\label{osc:H}
 \mathsf H^{X,Y}(a,v)
 =2^{d(v-a)}\int\Psi(r/2^v)
   \left(\int_{S^{d-1}}|U_X(r,w)|^2\,dw\right)
   \left(\int_{S^{d-1}}|V_Y(r,w)|^2\,dw\right)
   \frac{dr}{2^v}.
\end{equation}
It is defined without any separation condition.  In separated geometry
we also use the form with one shared angle
\begin{equation}\label{osc:HF}
 \mathsf H_F^{X,Y}(a,v)
 =2^{d(v-a)}\int\Psi(r/2^v)
   \left|\int_{S^{d-1}}\psi(w)
           U_X(r,w)\overline{V_Y(r,w)}\,dw\right|^2
   \frac{dr}{2^v}.
\end{equation}
These are values for the specified class laws, cells and endpoint
symbols, including the lists present in those symbols.  We use maxima over
finite collections of cell pairs and suprema when symbols vary; the family
in a supremum will be stated.
\end{definition}

Cauchy--Schwarz in $w$, with $0\le\psi\le1$, gives at each radius
\[
 \left|\int\psi(w)U_X(r,w)\overline{V_Y(r,w)}\,dw\right|^2
 \le \left(\int|U_X(r,w)|^2\,dw\right)
      \left(\int|V_Y(r,w)|^2\,dw\right).
\]
The radial measure and prefactor in the two forms are the same, so
\begin{equation}\label{osc:angular-cs}
 \mathsf H_F^{X,Y}(a,v)\le\mathsf H^{X,Y}(a,v)
\end{equation}
with exactly the same endpoint symbols and their support conditions.

There is a second operation available only in the positive independent
energy.  On a depth-$a$ cell, factor the reached mask factors out of its
transform as an angular function.  Lemma~\ref{osc:symbol-bounds} bounds
its modulus, so bounding that modulus inside the corresponding positive
angular integral gives, at bounded cost, the independent energy with
those factors deleted.  This operation keeps the remaining endpoint
symbols unchanged.  It does not permit deletion inside the shared
angular integral in $\mathsf H_F$ before its absolute square is taken.

\subsection{The weighted shell estimate}

The scalar input below includes a distance weight that cancels the leading
stationary distance power.  The conclusion retains the specified endpoint
lists and the same cell depth.

\begin{proposition}[Weighted shell conversion]\label{osc:shell-conversion}
Let $0\le a\le v\le N$ be integer depths, and let $X,Y$ be separated
cells at depth $a$ of one or two classes.  Let
$m_x(w)$ and $n_y(w)$ be products of the undifferentiated masks at one
index $i<J_0$ for two specified indexed descriptor lists.  Assume both
lists have at most $N^{C_{\mathcal A}}$ occurrences, where one exponent
$C_{\mathcal A}$ is fixed independently of $N$ for the whole family under
consideration, as in Definition~\ref{osc:symbols}.  Let $L$ be the largest
nominal length in either list, with $L=0$ when both lists are empty.
For measurable $|f|,|g|\le1$
define a possibly complex measure $\nu$ on the distance line by
\begin{equation}\label{osc:weighted-measure}
 \begin{aligned}
 \int F(s)\,d\nu(s)
 &=\frac1{M_1(a)M_2(a)}
   \int_{X\times Y} F(D(x,y))\,
       (D(x,y)2^a)^{-n_*/2}\\
 &\qquad\qquad\cdot f(x)\overline{g(y)}\,
       m_x(e(x,y))n_y(e(x,y))\,
       d\rho_1(x)d\rho_2(y).
 \end{aligned}
\end{equation}
If
\begin{equation}\label{osc:shell-regime}
 v-a\ge20\sigma N,\qquad L\le(v-a)/2+\sigma N,
\end{equation}
then for every fixed $T>0$,
\begin{equation}\label{osc:shell-estimate}
 2^{-a}\int_{2^{v-1}}^{2^{v+1}}|\widehat\nu(r)|^2\,dr
 \le R^{o(1)}\sup\mathsf H_F(a,v)+O(R^{-T}).
\end{equation}
The supremum has the same cells, class normalizations, lists and
smoothing index as in Equation~\eqref{osc:weighted-measure}.  Its symbols
are admissible with a derivative budget and common bounds for smooth angular factors
depending only on $T,d,C_{\mathcal A}$ and the fixed parameters.  All
these bounds are independent of $N$.
When $f=g=1$, the measure $\nu$ is positive.  The assertion also holds on
the fixed outer supports with $a=0$.
\end{proposition}

The rest of this Section proves Proposition~\ref{osc:shell-conversion}.
We first obtain the global sphere expansions, then separate the remaining
smooth spatial coefficient into bounded functions of the two endpoints.

\subsection{Stationary phase with global angular coefficients}

Give the sphere Laplacian the convention
$\Delta_{S^{n_*}}=\operatorname{div}\nabla$, so its eigenvalues are
nonpositive.  Define
\begin{equation}\label{osc:L-operators}
 \begin{split}
 c_{n_*}&=(2\pi)^{n_*/2}e^{-i\pi n_*/4},\\
 L_0&=c_{n_*}\operatorname{Id},\\
 L_j&=\frac{c_{n_*}}{(2i)^j j!}
   \prod_{q=0}^{j-1}
   \left(\Delta_{S^{n_*}}+q(q+1)
                    -\frac{n_*(n_*-2)}4\right),\qquad j\ge1.
 \end{split}
\end{equation}
For the opposite stationary point put
$L_{j,-1}b=\overline{L_j\overline b}$, and write $L_{j,+1}=L_j$.
These operators are defined globally, including when $d=2$.

\begin{lemma}[Direct and inverse sphere expansions]
\label{osc:stationary}
Suppose $b\in C^\infty(S^{n_*})$ obeys
$\|b\|_{C^m}\le C_m A^m$ at the finitely many orders needed below, where
$A\ge1$.  Let $\Delta\ne0$, $e=\Delta/|\Delta|$,
$\Lambda=r|\Delta|$, and $r>0$.  A cutoff $\psi$ is assumed to equal one
on a fixed neighborhood of $e$ and zero on a fixed neighborhood of
$-e$, with uniformly bounded derivatives.  For every integer $J\ge1$,
\begin{equation}\label{osc:direct}
 \begin{aligned}
 \int_{S^{n_*}}e^{irw\cdot\Delta}\psi(w)b(w)\,dw
 &=\Lambda^{-n_*/2}
   \left[e^{i\Lambda}\sum_{j=0}^{J-1}
                    \Lambda^{-j}(L_jb)(e)+E_J\right],\\
 |E_J|&\le C_J\Lambda^{-J}A^{2J+d+2},
 \end{aligned}
\end{equation}
provided $\Lambda\ge1$.  Without a cutoff the corresponding formula is
\begin{equation}\label{osc:full-sphere}
 \int_{S^{n_*}}e^{irw\cdot\Delta}b(w)\,dw
 =\Lambda^{-n_*/2}
   \left[\sum_{\epsilon\in\{+1,-1\}}e^{i\epsilon\Lambda}
       \sum_{j=0}^{J-1}\Lambda^{-j}
                 (L_{j,\epsilon}b)(\epsilon e)+E_J'\right],
\end{equation}
with the same remainder bound.  Define global differential operators
\begin{equation}\label{osc:B-operators}
 B_0=L_0^{-1},\qquad
 B_j=-L_0^{-1}\sum_{h=1}^j L_h B_{j-h}.
\end{equation}
Each $B_j$ is a polynomial in $\Delta_{S^{n_*}}$ of differential order at
most $2j$.  If, for a fixed $c>0$,
\begin{equation}\label{osc:stationary-regime}
 \Lambda\ge R^{c\sigma},\qquad
 A^2/\Lambda\le R^{-c\sigma},\qquad 1\le A\le R^2,
\end{equation}
then for every prescribed $T>0$ a fixed sufficiently large $J$ gives
\begin{equation}\label{osc:inverse}
 e^{i\Lambda}b(e)
 =\Lambda^{n_*/2}\sum_{j=0}^{J-1}\Lambda^{-j}
       \int_{S^{n_*}}e^{irw\cdot\Delta}\psi(w)(B_jb)(w)\,dw
       +O(R^{-T}).
\end{equation}
The extracted remainders in Equations~\eqref{osc:direct} and
\eqref{osc:full-sphere} are also $O(R^{-T})$ after increasing $J$.
All constants are uniform in $e$ in the stipulated geometry.

For $r=2^v s$ with $s$ in a fixed compact subinterval of $(0,\infty)$,
each coefficient $\Lambda^{-j}(L_{j,\epsilon}b)(\epsilon e)$ has uniformly
bounded derivatives of every fixed order in $s$ under
Equation~\eqref{osc:stationary-regime}.  Every differential coefficient
retains the supports of all the mask factors in $b$.
\end{lemma}

\begin{proof}
We first obtain a local expansion with a controlled remainder.  Away
from $e$ and $-e$, the angular gradient of $w\cdot e$ is bounded below.
Repeated integration by parts with this gradient bounds the integral
there by any chosen power of $\Lambda^{-1}$, paying at most the
corresponding angular derivatives of $b$.  Near the maximum, rotate $e$
to the last coordinate vector and use
\[
 w(u)=\bigl(u\sqrt{1-|u|^2/4},1-|u|^2/2\bigr),\qquad
 dw=(1-|u|^2/4)^{(n_*-2)/2}\,du.
\]
The phase is exactly $w\cdot e=1-|u|^2/2$ and the Jacobian equals one at
$u=0$.  These charts and their
derivatives can be chosen uniformly on finitely many patches of possible
$e$.  Thus the local integral is
$e^{i\Lambda}\int e^{-i\Lambda|u|^2/2}a_e(u)\,du$, with $a_e$
compactly supported in a fixed coordinate ball and
$\|a_e\|_{C^m}\lesssim_m A^m$.

The Fourier formula for a Gaussian, obtained first with positive damping
and then by taking its limit, expresses the quadratic integral
$\int e^{-i\Lambda|u|^2/2}a_e(u)\,du$ as
$c_{n_*}\Lambda^{-n_*/2}$ times the integral of
$\widehat a_e(\xi)e^{i|\xi|^2/(2\Lambda)}$ with the Fourier inversion
normalization.  Taylor expansion through order $J-1$ has remainder at
most
\[
 C_J\Lambda^{-J}
       \int |\xi|^{2J}|\widehat a_e(\xi)|\,d\xi
 \le C_J\Lambda^{-J}A^{2J+d+2}.
\]
The last inequality follows by integrating by parts in the transform to
an order greater than $2J+n_*$; all supports have fixed diameter.
For the nonstationary portion take
$M=J+\lceil n_*/2\rceil$ integrations by parts.  After extracting
$\Lambda^{-n_*/2}$, its bound is
$C_M\Lambda^{n_*/2-M}A^M\le C_J\Lambda^{-J}A^{2J+d+2}$.
We have proved
Equation~\eqref{osc:direct} with differential coefficients of order at
most $2j$ and leading coefficient $c_{n_*}b(e)$.

We identify these coefficients without making a choice of moving
endpoint coordinates.  For this identification fix a smooth $b$ and use
a polar cutoff $\chi(w\cdot e)$ equal to one near $w=e$ and zero outside
a slightly larger fixed cap.  The Fourier extension
$F(\xi)=\int e^{iw\cdot\xi}b(w)\,dw$ satisfies
$(\Delta_\xi+1)F=0$.  The same identity for its polar-cutoff outgoing
part has an error decreasing faster than every power of $|\xi|$:
derivatives of the cutoff are supported away from both stationary
points, so the preceding integration by parts applies, also after any
fixed number of derivatives.  In polar coordinates the differential
operator is
\[
 \partial_\Lambda^2+\frac{n_*}{\Lambda}\partial_\Lambda
       +1+\Lambda^{-2}\Delta_{S^{n_*}}.
\]
The Gaussian expansion can be differentiated any fixed number of times
in its angular and radial parameters: differentiate the compactly
supported coordinate amplitudes and Gaussian formula, and take more
terms to bound the differentiated remainder.  We may therefore
substitute the local outgoing expansion
$e^{i\Lambda}\Lambda^{-n_*/2}
\sum_{j\ge0}\Lambda^{-j}a_j(e)$.  Uniqueness of its power coefficients
gives
\[
 2i(j+1)a_{j+1}
  =\left(\Delta_{S^{n_*}}+j(j+1)
                     -\frac{n_*(n_*-2)}4\right)a_j.
\]
Starting from $a_0=c_{n_*}b$, this is exactly
Equation~\eqref{osc:L-operators}.  A cutoff equal to one near $e$ has the
same coefficients, since the difference is nonstationary.  The minimum
has the conjugate Gaussian phase, proving Equation~\eqref{osc:full-sphere}.

For inversion, apply Equation~\eqref{osc:direct} to $B_jb$, expanding its
series through index $J-j-1$.  Equation~\eqref{osc:B-operators} cancels
every coefficient of total index $1,\ldots,J-1$, and its index-zero
coefficient is $b$.  The derivative bound for $B_jb$ at order $m$ is
$C_{j,m}A^{m+2j}$.  After multiplication by $\Lambda^{-j}$ its extracted
remainder is therefore at most
$C_J\Lambda^{-J}A^{2J+d+2}$, uniformly in $0\le j<J$.
Under Equation~\eqref{osc:stationary-regime} this is at most
\[
 C_J(A^2/\Lambda)^J A^{d+2}
 \le C_J R^{-c\sigma J+2(d+2)}.
\]
Taking $J$ large proves all the asserted arbitrary-power estimates.

Finally, $e$ and $b$ are independent of $s$, and each $s$ derivative of
$\Lambda^{-j}(L_{j,\epsilon}b)(\epsilon e)$ costs only a fixed constant
times $\Lambda^{-j}A^{2j}$, which is bounded.  The same elementary
Taylor remainder argument can accommodate any prescribed finite number
of radial derivatives by increasing the truncation order.  In the
present applications $\Lambda\lesssim R^2$, so even unstripped radial
derivatives of the error can be made an arbitrary negative power by
increasing that order once more.  For the support assertion, use
\[
 \Delta_{S^{d-1}}=\sum_{u<v}\Omega_{uv}^2.
\]
Leibniz expansion of any $L_j$ or $B_j$ differentiates, but never removes,
each original mask factor.  Derivatives of a smooth function cannot
enlarge its support.  This proves the assertion, including for arbitrary
measurable spatial multipliers.
\end{proof}

\subsection{Separation of the remaining spatial coefficient}

\begin{lemma}[Separated spatial expansion]\label{osc:spatial-separation}
On separated depth-$a$ cells, for every fixed $j\ge0$ there are bounded
measurable functions $f_\ell$ on $X$ and $g_\ell$ on $Y$ and constants
$c_\ell$ such that
\begin{equation}\label{osc:spatial-series}
 (D(x,y)2^a)^{-j}
   =\sum_\ell c_\ell f_\ell(x)g_\ell(y),\qquad
 |f_\ell|,|g_\ell|\le1,\qquad
 \sum_\ell|c_\ell|\le C_{d,j}.
\end{equation}
The convergence is uniform.  The same assertion holds in the fixed
outer geometry, with a constant depending on that geometry.
\end{lemma}

\begin{proof}
Write $x=c_X+2^{-a}u$, $y=c_Y+2^{-a}z$.  Both rescaled endpoint cells
lie in fixed boxes.  By Equation~\eqref{osc:separation},
$|2^a(c_X-c_Y)+u-z|$ is bounded above and below by positive constants on
a fixed neighborhood of their product.  Consequently its negative
$j$th power has a smooth extension, with uniform bounds of every fixed
order, supported in a slightly larger product box.  Regard this
extension as a smooth function on a fixed torus of dimension $2d$.
Its Fourier series has coefficients bounded by
$C_m(1+|k|+|l|)^{-m}$.  Choosing $m>2d$ makes their absolute sum bounded.
Each Fourier exponential splits into one function of $u$ and one of
$z$, each of modulus one.  Pulling them back proves the claim.
The proof on the fixed separated outer supports is identical.
\end{proof}

\subsection{Proof of the weighted shell estimate}

\begin{proof}[Proof of Proposition~\ref{osc:shell-conversion}]
Put $q=v-a$, $s=r/2^v$, and
$D_{\max}=\max(1,2^L R^{-50\sigma})$.  On the shell in question,
$\Lambda=rD(x,y)\asymp2^q$ and
\[
 D_*=R^\sigma D_{\max}
 \le\max(R^\sigma,2^{q/2}R^{-48\sigma}),\qquad
 \frac{D_*^2}{\Lambda}
 \lesssim\max(R^{2\sigma}2^{-q},R^{-96\sigma})
 \lesssim R^{-18\sigma}.
\]
Moreover $\Lambda\gtrsim R^{20\sigma}$ and $D_*\le R^2$.
Fixed geometric constants can be absorbed by slightly decreasing these
positive exponents.  Lemma~\ref{osc:stationary} therefore applies
uniformly to the endpoint product
$m_x(w)n_y(w)$.

We first estimate $\widehat\nu(-r)$, which has the outgoing phase
$e^{irD}$.  This suffices also for $\widehat\nu(r)$: conjugating the
spatial multipliers conjugates $\nu$, and
$\widehat{\overline\nu}(-r)=\overline{\widehat\nu(r)}$.
Use Equation~\eqref{osc:inverse} pointwise for each $(x,y)$ and multiply
by the weight in Equation~\eqref{osc:weighted-measure}.  The leading
distance power cancels exactly:
\begin{equation}\label{osc:weight-cancellation}
 (D2^a)^{-n_*/2}(rD)^{n_*/2}(rD)^{-j}
 = (r2^{-a})^{n_*/2}(r2^{-a})^{-j}(D2^a)^{-j}.
\end{equation}
The operator $B_j$ is global.  Its Leibniz expansion on $m_xn_y$ is a
sum of products of angular derivatives of the individual masks, of
total order at most $2j$.  Divide every differentiated mask by its own
scale $D_{l_{\mathfrak t}}$ to the corresponding order.  The resulting factors are
admissible, retain every descriptor, and their extracted scales have
product at most $D_{\max}^{2j}$.  Their number is polynomial in $N$,
since $j$ is fixed.  Smooth angular factors, if produced, can be
included in either endpoint symbol.

The inverse factor pays for all these scales, since
\begin{equation}\label{osc:scale-payment}
 2^{-jq}D_{\max}^{2j}
 \le\bigl[\max(2^{-q},R^{-98\sigma})\bigr]^j\le1.
\end{equation}
Apply Lemma~\ref{osc:spatial-separation} to the remaining factor
$(D2^a)^{-j}$.  Its absolutely summable expansion only multiplies the
original arbitrary spatial factors by bounded endpoint functions.
After summing the finitely many inverse orders, we obtain
\begin{equation}\label{osc:shell-factorization}
 \widehat\nu(-2^vs)
  =2^{n_*q/2}\sum_\alpha c_\alpha(s)
       \int\psi(w)U_{X,\alpha}(2^vs,w)
                       \overline{V_{Y,\alpha}(2^vs,w)}\,dw
       +\mathcal E(2^vs),
\end{equation}
where
\[
 \sum_\alpha\|c_\alpha\|_{L^\infty([1/2,2])}
 \le R^{o(1)}.
\]
The factors $s^{n_*/2-j}$ are included in $c_\alpha(s)$ and are
uniformly bounded on this interval.  The index set can be countable
because of the spatial Fourier series; its absolute summability
justifies all interchanges.  All endpoint transforms in
Equation~\eqref{osc:shell-factorization} have the form
Equation~\eqref{osc:transforms} with the stated enlarged fixed budget.

The pointwise error $\mathcal E$ is smaller than any prescribed negative
power of $R$.  Indeed Lemma~\ref{osc:stationary} supplies such an error
before the endpoint integrations.  The rescaled distance weight is
bounded, and the normalized cell masses are
$\rho_j(X)/M_j(a)\le R^{O_d(\sigma)}$.  Increasing the initial accuracy
absorbs these factors and any subsequent square.  The shell length is at
most $2^{N+1}=2R$, since $v\le N$, and is absorbed by the same choice.

Apply the triangle inequality in $L^2([1/2,2],ds)$ to
Equation~\eqref{osc:shell-factorization}.  Since $\Psi\ge1$ there, the
square of each angular integral is bounded by its contribution to the
corresponding $\mathsf H_F$ energy.  The complete prefactor is
\begin{equation}\label{osc:normalization-identity}
 2^{-a}\,2^v\,2^{n_*(v-a)}=2^{d(v-a)}.
\end{equation}
This is precisely the normalization in Equation~\eqref{osc:HF}, proving
Equation~\eqref{osc:shell-estimate}.  Complex conjugation gives the
positive-frequency formulation stated in the Proposition.
\end{proof}
